\subsection{A singleton adjacent arrow and a binary composition map}

\begin{proposition}[Proved: the singleton--binary triangle for arbitrary corners]
\label{CP2:A-singleton-binary-triangle-all-corners}
Let $C$ be a three-object category with full linear acyclic support.
If
\[
 |C(0,1)|=1,\qquad |C(1,2)|=|C(0,2)|=2,
\]
then all rational degree-two products of $BC$ vanish.  There is
no restriction on the three finite corner monoids.  The opposite
statement holds as well.
\end{proposition}
\begin{proof}
Write $M=C(0,0)$, $B=C(1,1)$, $N=C(2,2)$ and
$C(0,1)=\{f\}$.  Composition defines an $N$-equivariant map
\[
 \lambda:C(1,2)\longrightarrow C(0,2),\qquad g\longmapsto gf.
\]
Since both sets have two elements, this map is either constant or
bijective.

Suppose it is constant, with value $z$.  For every $n\in N$,
$nz=n\lambda(g)=\lambda(ng)=z$, so $z$ is an $N$-fixed
connecting arrow.  The full block $D=C|_{\{0,1\}}$ has its
connecting hom-set a singleton; its classifying space is
$BM*BB$.  Both factors are connected, hence this join has
$H^2(BD;\mathbb Q)=0$.  Regard $C$ as the collage $D\to N$.
The fixed arrow $z$, with the identity string at its source object,
gives a section from $BN$ into the relevant action-category
subspace of the bridge.  Consequently the bridge pullback from
$H^2(BN;\mathbb Q)$ is injective.  Compatibility of the two
corner restrictions of a global second class, and the zero
restriction to $BD$, force its restriction to $BN$ to be zero.
The two-block relative-suspension argument then kills all products
of second classes.

Suppose instead that $\lambda$ is bijective.  For every $m\in M$,
the unique arrow $f$ satisfies $fm=f$.  Thus
\[
       \lambda(g)m=(gf)m=g(fm)=gf=\lambda(g).
\]
Surjectivity of $\lambda$ proves that the entire right $M$ action
on $C(0,2)$ is trivial.  Its action on the singleton $C(0,1)$
is also trivial.  The actual retraction replacing $M$ by its
identity therefore induces an isomorphism in rational second
cohomology, by Lemma~\ref{CP2:A-trivial-cut-H2}, or equivalently
Theorem~\ref{CP2:A-trivial-extreme-corner-retract}.
Denote the resulting category by $C'$.

Delete the middle object of $C'$.  For its inclusion $j$ of the
full deletion, the only object of $(j\downarrow1)$ is the unique
arrow from the now trivial initial corner to object $1$, and its
only endomorphism is the identity.  At either retained object the
comma category has its usual terminal object.  Quillen's
Theorem~A therefore makes this deletion a homotopy equivalence.
The resulting two-object category is the ordinary left-$N$ act
cone on the two-element set $C(0,2)$.  Its rational degree-two
products vanish by Lemma~\ref{CP2:A-high-affine}.

Every second class of $BC$ is the pullback of a second class of
$BC'$ through the actual retraction, and pullback preserves cup
products.  This proves the claim in the bijective case also.
Taking opposites gives the final assertion.
\end{proof}
